Assume now that the weights are continuous. If the obstacle is imposed only quasi-everywhere on $K$, the regularized envelope itself becomes a competitor, so the supremum needs no regularization; a constraint imposed at every point would miss the values created by regularization. Indeed, for $q\in C(K,\mathbb R)$,
\begin{equation}\label{eq:affine-qe-pointwise-envelope}
 F_{\delta,q}(z)=\sup\{u(z):u\in\mathcal L_\delta,
                          \ u\leq q\text{ q.e. on }K\}.
\end{equation}
The envelope $F_{\delta,q}$ itself is a competitor on the right, since it agrees with the raw envelope off a pluripolar set by \cref{prop:affine-projective-envelope}. Conversely, the argument of \cref{cor:affine-polar-obstacle-invariance} shows that every competitor on the right lies below $F_{\delta,q}$ everywhere.

For $a\in\mathbb C^n$, let $\mathcal J_a(K)$ be the set of \emph{supporting pairs} $(\mu,c)$, where $\mu$ is a Borel probability measure on $K$ and $c\geq0$, such that
\begin{equation}\label{eq:affine-supporting-probability-tests}
 u(a)\leq\int_K f\,\mathrm{d}\mu+c
\end{equation}
whenever $u\in\mathcal L_1$, $f\in C(K,\mathbb R)$ and $u\leq f$ quasi-everywhere on $K$. We call $\mu$ a \emph{supporting probability} at $a$.

\begin{theorem}\label{thm:affine-supporting-probabilities}
For every $a\in\mathbb C^n$, $q\in C(K,\mathbb R)$ and $\delta>0$,
\begin{equation}\label{eq:affine-supporting-probability-minimum}
 F_{\delta,q}(a)=\min_{(\mu,c)\in\mathcal J_a(K)}
             \left(\int_K q\,\mathrm{d}\mu+\delta c\right).
\end{equation}
In particular, $\mathcal J_a(K)$ is nonempty and
\[
 \min_{(\mu,c)\in\mathcal J_a(K)}c=V_K^*(a).
\]
\end{theorem}
\begin{proof}
On the real Banach space $C(K)$, put $T(f)=F_{1,f}(a)$. By \eqref{eq:affine-qe-pointwise-envelope}, $T$ is concave, since convex combinations of competitors are competitors. Comparison of obstacles and translation by constants show that $T$ is nondecreasing and satisfies
\[
 T(f+t)=T(f)+t,\qquad |T(f)-T(g)|\leq\|f-g\|_\infty.
\]
Fix $f_0\in C(K)$. The directional derivative of the convex Lipschitz function $-T$ at $f_0$ is finite and sublinear. By the Hahn--Banach theorem, there is a continuous linear functional dominated by it, and since difference quotients of a convex function are monotone, this functional supports $-T$ at $f_0$. Changing the sign, we obtain a continuous linear functional $\ell$ with
\[
 T(f)\leq T(f_0)+\ell(f-f_0)\qquad(f\in C(K)).
\]
Testing $f=f_0+th$ with $h\geq0$ and $t>0$, and using monotonicity, gives $\ell(h)\geq0$. Testing with $f=f_0\pm t$ gives $\ell(1)=1$. By the Riesz representation theorem, $\ell(f)=\int_K f\,\mathrm{d}\mu$ for a Borel probability measure $\mu$ on $K$. Set $c=T(f_0)-\int_K f_0\,\mathrm{d}\mu$. The supporting inequality at $f=0$ gives $c\geq T(0)=V_K^*(a)\geq0$. Together with \eqref{eq:affine-qe-pointwise-envelope}, it shows that $(\mu,c)\in\mathcal J_a(K)$ and
\[
 T(f_0)=\int_K f_0\,\mathrm{d}\mu+c.
\]
Conversely, applying \eqref{eq:affine-supporting-probability-tests} to $u=F_{1,f_0}$, which lies below $f_0$ quasi-everywhere on $K$, shows that $T(f_0)\leq\int_Kf_0\,\mathrm{d}\mu+c$ for every pair in $\mathcal J_a(K)$. This proves \eqref{eq:affine-supporting-probability-minimum} for $\delta=1$, with the minimum attained. The identity $F_{\delta,q}=\delta F_{1,q/\delta}$ gives the general case, and the zero obstacle gives the last assertion.
\end{proof}
\index{Envelope!supporting probability}

\begin{theorem}\label{thm:affine-supported-growth-derivative}
If $\psi\in C(K,\mathbb R)$, then for every $a\in\mathbb C^n$ and $\tau>0$,
\begin{equation}\label{eq:affine-supported-growth-derivative}
 D_\tau(a)=\min\left\{c:(\mu,c)\in\mathcal J_a(K),
                \ \int_K\psi\,\mathrm{d}\mu+\tau c=F_\tau(a)\right\}.
\end{equation}
Consequently, $D_\tau(a)=V_K^*(a)$ if and only if there is a Borel probability measure $\mu$ on $K$ such that
\[
 (\mu,V_K^*(a))\in\mathcal J_a(K),\qquad
 \int_K\psi\,\mathrm{d}\mu=F_\tau(a)-\tau V_K^*(a).
\]
\end{theorem}
\begin{proof}
Each defining inequality of $\mathcal J_a(K)$ is preserved under weak convergence of the probability measures and convergence of the constants, because the test function $f$ is continuous. Since the Borel probability measures on the compact metric space $K$ form a weakly compact set, the pairs in $\mathcal J_a(K)$ with $c\leq C$ form a compact set for each finite $C$.

Call a pair in $\mathcal J_a(K)$ \emph{active} at $\tau$ if it satisfies the equality in \eqref{eq:affine-supported-growth-derivative}. Such pairs exist by \cref{thm:affine-supporting-probabilities}, and since $\int_K\psi\,\mathrm{d}\mu\geq m=\inf_K\psi$, they satisfy
\[
 0\leq c\leq\frac{F_\tau(a)-m}{\tau}.
\]
The set of active pairs is therefore compact, and the minimum of $c$ over it is attained. By \cref{thm:affine-supporting-probabilities} at $\tau+h$, each active pair gives $F_{\tau+h}(a)\leq F_\tau(a)+hc$, hence $D_\tau(a)\leq c$.

Conversely, for $h>0$ choose a pair $(\mu_h,c_h)$ attaining the minimum in \eqref{eq:affine-supporting-probability-minimum} at $\tau+h$, with $q=\psi$. The representation at $\tau$ gives
\[
 F_\tau(a)\leq\int_K\psi\,\mathrm{d}\mu_h+\tau c_h
            =F_{\tau+h}(a)-hc_h.
\]
Thus $0\leq c_h\leq(F_{\tau+h}(a)-F_\tau(a))/h$, and these constants are bounded by \eqref{eq:affine-initial-slope-bounds}. Along a sequence $h\downarrow0$, compactness gives a limiting pair $(\mu,c)\in\mathcal J_a(K)$. The same bound makes $F_{\tau+h}(a)\to F_\tau(a)$ and $hc_h\to0$, and continuity of $\psi$ then gives $\int_K\psi\,\mathrm{d}\mu+\tau c=F_\tau(a)$. Thus the limit pair is active, and it satisfies $c\leq D_\tau(a)$. This proves \eqref{eq:affine-supported-growth-derivative}. Finally, every pair has $c\geq V_K^*(a)$ by \cref{thm:affine-supporting-probabilities}, so the attainment of the minimum gives the stated criterion.
\end{proof}

At a point $a$ of regular contact, $F_\tau(a)=\psi(a)$, so the criterion involves only the original weight and tests by functions in $\mathcal L_1$. If the defining tests instead required $u\leq f$ at every point of $K$, the pair formed by the unit point mass at $a\in K$ and the constant $0$ would always be allowed. The resulting infimum would then be at most $\psi(a)$, so it could not detect a defect created by regularization.

\subsubsection{Variation of continuous weights}
\label{sec:affine-weight-response}

Supporting probabilities describe changes of the obstacle as well as the growth allowance. Keep the non-pluripolar compact $K\subset\mathbb C^n$ fixed. For $a\in\mathbb C^n$, $\delta>0$ and $q\in C(K,\mathbb R)$, put
\begin{equation}\label{eq:affine-active-supports}
 \mathcal A_a(\delta,q)=
 \left\{(\mu,c)\in\mathcal J_a(K):
   \int_Kq\,\mathrm{d}\mu+\delta c=F_{\delta,q}(a)\right\}.
\end{equation}
The weights carry the uniform norm, probabilities carry the weak topology, and constants carry the ordinary topology.

\begin{lemma}\label{lem:affine-active-support-stability}
The set $\mathcal A_a(\delta,q)$ is nonempty, convex and compact. For $\delta,\eta\geq\epsilon>0$ and $p,q\in C(K,\mathbb R)$,
\begin{equation}\label{eq:affine-weight-growth-lipschitz}
 |F_{\eta,p}(a)-F_{\delta,q}(a)|
 \leq\|p-q\|_\infty+
 |\eta-\delta|\left(V_K^*(a)+
                    \frac{\operatorname{osc}_Kq}{\epsilon}\right).
\end{equation}
If $\delta_j\to\delta>0$, $q_j\to q$ uniformly and $(\mu_j,c_j)\in\mathcal A_a(\delta_j,q_j)$, then these pairs have a convergent subsequence, and every subsequential limit belongs to $\mathcal A_a(\delta,q)$.
\end{lemma}
\begin{proof}
Nonemptiness follows from \cref{thm:affine-supporting-probabilities}. The continuous test inequalities make $\mathcal J_a(K)$ closed and convex. The active equality is also closed and affine. For an active pair, \eqref{eq:affine-obstacle-sandwich} gives
\[
 0\leq c\leq\frac{F_{\delta,q}(a)-\inf_Kq}{\delta}
 \leq V_K^*(a)+\frac{\operatorname{osc}_Kq}{\delta}.
\]
The bounded constant slice is compact, as in the proof of \cref{thm:affine-supported-growth-derivative}; hence the active set is compact.

With $q$ fixed, \cref{lem:affine-growth-increment-bounds} and \eqref{eq:affine-obstacle-sandwich} bound a growth increment by its length times $V_K^*(a)+\operatorname{osc}_Kq/\epsilon$. Obstacle order and constant translation bound the change from $q$ to $p$ at fixed $\eta$ by $\|p-q\|_\infty$. Their sum is \eqref{eq:affine-weight-growth-lipschitz}.

For the sequence in the statement, the active constants have a common bound because the weights are uniformly bounded and the growth allowances stay bounded below. Weak compactness of probabilities gives subsequential pair convergence. Closedness retains membership in $\mathcal J_a(K)$. The displayed Lipschitz estimate gives envelope-value convergence, while uniform convergence of $q_j$ and weak convergence of $\mu_j$ give convergence of the integrals. Passing to the active equality proves the last assertion.
\end{proof}
\index{Envelope!weight variation}

\begin{proposition}\label{prop:affine-joint-directional-response}
For $r\in C(K,\mathbb R)$ and $s\in\mathbb R$, one has
\begin{equation}\label{eq:affine-joint-directional-response}
 \lim_{h\downarrow0}
 \frac{F_{\delta+hs,q+hr}(a)-F_{\delta,q}(a)}h
 =\min_{(\mu,c)\in\mathcal A_a(\delta,q)}
       \left(\int_Kr\,\mathrm{d}\mu+sc\right).
\end{equation}
The same limit holds along $h_j\downarrow0$ with $s_j\to s$ and $r_j\to r$ uniformly, replacing $s,r$ in the quotient by $s_j,r_j$. Thus the one-sided directional derivative is stable under convergent directions.
\end{proposition}
\begin{proof}
Compactness of the active set attains the minimum on the right. For a fixed base active pair, the representation at the perturbed parameters gives
\[
 F_{\delta+hs,q+hr}(a)
 \leq F_{\delta,q}(a)+h\left(\int_Kr\,\mathrm{d}\mu+sc\right).
\]
A minimizing base pair therefore bounds the limsup by the stated minimum. Conversely, choose a minimizing pair $(\mu_h,c_h)$ at $(\delta+hs,q+hr)$. The representation at the base parameters gives
\[
 \frac{F_{\delta+hs,q+hr}(a)-F_{\delta,q}(a)}h
 \geq\int_Kr\,\mathrm{d}\mu_h+sc_h.
\]
The local Lipschitz estimate bounds the quotients. Along a sequence approaching their liminf, \cref{lem:affine-active-support-stability} supplies a convergent subsequence of pairs with a base active limit. Its directional value is at least the minimum. This proves \eqref{eq:affine-joint-directional-response}.

For convergent directions, the same two inequalities hold. The fixed minimizing base pair gives a convergent upper bound. Perturbed minimizing pairs still have uniformly bounded constants and base active limits. Uniform convergence of $r_j$ controls their moving integrals, and $s_j\to s$ controls the constant terms. The lower bound therefore has the same limit. All perturbed growth allowances are positive for sufficiently small increments.
\end{proof}

\begin{corollary}\label{cor:affine-two-sided-growth-slopes}
For fixed $a$ and $q$, the function $\delta\mapsto F_{\delta,q}(a)$ has finite left and right derivatives $D^-_{\delta,q}(a)$ and $D^+_{\delta,q}(a)$ at every positive parameter, with
\begin{equation}\label{eq:affine-two-sided-growth-slopes}
 D^-_{\delta,q}(a)=\max_{\mathcal A_a(\delta,q)}c,
 \qquad
 D^+_{\delta,q}(a)=\min_{\mathcal A_a(\delta,q)}c.
\end{equation}
It is differentiable exactly when all active pairs have the same constant $c$. For each fixed $a,q$, the set of positive parameters where it fails to be differentiable is at most countable.
\end{corollary}
\begin{proof}
Apply \eqref{eq:affine-joint-directional-response} with $r=0$ and $s=1,-1$. The negative direction gives minus the left derivative, proving \eqref{eq:affine-two-sided-growth-slopes}. Concavity gives
\[
 D^+_{\alpha,q}(a)\geq D^-_{\beta,q}(a)
 \qquad(0<\alpha<\beta).
\]
Hence the nonempty open intervals $(D^+_{\delta,q}(a),D^-_{\delta,q}(a))$ at exceptional parameters are pairwise disjoint. Assigning a rational number in each interval proves countability.
\end{proof}
\index{Envelope!growth differentiability}

\begin{corollary}\label{cor:affine-weight-differentiability}
At fixed $a$ and $\delta>0$, the map $q\mapsto F_{\delta,q}(a)$ on $C(K,\mathbb R)$ is G\^ateaux differentiable at $q$ if and only if the active probability is unique. In that case its derivative is
\[
 r\longmapsto\int_Kr\,\mathrm{d}\mu.
\]
\end{corollary}
\begin{proof}
With $s=0$, \eqref{eq:affine-joint-directional-response} gives the minimum of $\int_Kr\,\mathrm{d}\mu$ over active pairs. The negative direction gives minus the maximum. If the active probability is unique, these values give the same continuous linear two-sided derivative. Conversely, a linear two-sided derivative forces the minimum and maximum to agree for every continuous $r$. Such tests separate Borel probabilities on $K$, so all active probabilities coincide. Their constant is then fixed by \eqref{eq:affine-active-supports} and $\delta>0$.
\end{proof}
\index{Envelope!weight differentiability}

\begin{example}\label{ex:affine-quadratic-weight-corner}
Let $K=\overline D\subset\mathbb C$ be the closed unit disk and fix $\delta>0$. For $q_t(z)=-t|z|^2$,
\begin{equation}\label{eq:affine-quadratic-weight-corner}
 F_{\delta,q_t}(0)=
 \begin{cases}
 -t,&t\geq0,\\
 0,&-\delta/2\leq t\leq0.
 \end{cases}
\end{equation}
Thus the weight functional at the origin is not G\^ateaux differentiable at the zero weight, and it has more than one active probability there.
\end{example}
\begin{proof}
For $t\geq0$, each raw competitor is at most $-t$ on the boundary circle. The subharmonic maximum principle bounds it by $-t$ on $D$. Regularization retains that bound near zero, and the constant $-t$ is a global competitor.

For $t=-b$ with $0\leq b\leq\delta/2$, define
\[
 u_b(z)=
 \begin{cases}
 b|z|^2,&|z|\leq1,\\
 b+\delta\log|z|,&|z|>1.
 \end{cases}
\]
This function is continuous and subharmonic. Its ordinary Laplacian has interior density $4b$, zero exterior density and a nonnegative circle contribution from the outward derivative jump $\delta-2b$. It belongs to $\mathcal L_\delta$ and meets $q_{-b}$ on $K$. Hence $F_{\delta,q_{-b}}(0)\geq0$. On the open disk the raw envelope is at most the continuous function $b|z|^2$, so regularization gives the reverse inequality at zero. This proves \eqref{eq:affine-quadratic-weight-corner}. The right and left derivatives in $t$ are $-1$ and $0$. \Cref{cor:affine-weight-differentiability} gives the final assertion.
\end{proof}


\subsubsection{Local regularity and strict weights}
\label{sec:affine-local-regularity}

Local contact depends on the compact pieces near the point. The relative extremal function of \cref{def:relative-extremal-function} detects their unweighted regularity. A strict lower Hessian bound on a global extending weight then turns failure of that regularity into a quantitative weighted defect.

For a non-pluripolar compact $E\subset\mathbb C^n$, write $V_E$ and $V_E^*$ for its raw and regularized growth-one envelopes with zero obstacle. For other weights and growth allowances, write $F_{\delta,E,q}$ for the regularized envelope on $E$ with the everywhere obstacle $q$.

\begin{lemma}\label{lem:relative-affine-zero-value}
Let $0<r<R$, $B=B(a,R)$, and let $E\subseteq\overline B(a,r)$ be compact. If $E$ is pluripolar, then $h_{E,B}^*=0$ on $B$. If $E$ is non-pluripolar, then
\[
 h_{E,B}^*(a)=-1\quad\Longleftrightarrow\quad V_E^*(a)=0.
\]
\end{lemma}
\begin{proof}
In the pluripolar case, choose $b\in\mathcal L_1$ with poles on $E$ by \cref{lem:affine-logarithmic-pole-barrier}, and choose a finite upper bound $C$ for $b$ on $B$. The relative competitors
\[
 u_\epsilon=\max\{-1,\epsilon(b-C)\}
\]
tend to zero almost everywhere as $\epsilon\downarrow0$. Their upper bound $h_{E,B}^*$ is therefore nonnegative almost everywhere, then everywhere by canonical order. Since it is also nonpositive, it vanishes.

Suppose $E$ is non-pluripolar. \Cref{prop:affine-projective-envelope} gives a finite upper bound $M\geq1$ for $V_E^*$ on $B$. For each raw global competitor $v$, the function $\max(-1,v/M-1)$ is a relative competitor. Taking suprema and regularizing gives
\[
 h_{E,B}^*(a)\geq\max\{-1,V_E^*(a)/M-1\}.
\]
As $V_E^*\geq0$, the relative value $-1$ forces the global value zero.

For the converse choose $r<R_0<R$ and $k>0$ with $k\log(R_0/r)>1$. The plurisubharmonic function
\[
 L(z)=k\log\bigl(\lVert z-a\rVert/r\bigr)
\]
has growth allowance $k$, with value $-\infty$ at $a$. For a relative competitor $u$, paste $W=\max(u+1,L)$ on $B$ with $W=L$ outside $B$. On the outer collar $L>1\geq u+1$, so the pieces agree and \cref{lem:qpsh-gluing} applies. On $E$ one has $W=0$; on the open inner ball one has $W=u+1$. Hence $W/k$ is a global growth-one competitor, and the same $k$ works for the whole relative family. Supremum and regularization near $a$ give
\[
 V_E^*(a)\geq(1+h_{E,B}^*(a))/k.
\]
A zero global value thus forces the relative value $-1$.
\end{proof}

Fix again a non-pluripolar compact $K\subset\mathbb C^n$. For $a\in K$ put $K_r=K\cap\overline B(a,r)$.

\begin{definition}\label{def:affine-local-pluriregularity}\index{Pluriregular compact!local}
The compact $K$ is locally pluriregular at $a$ if every $K_r$, $r>0$, is non-pluripolar and $V_{K_r}^*(a)=0$. It is locally $(\delta,q)$-regular there if $F_{\delta,K_r,q}(a)=q(a)$ for every $r>0$.
\end{definition}

It suffices to check arbitrarily small radii: compact-set monotonicity handles the larger pieces. \Cref{lem:relative-affine-zero-value} identifies local pluriregularity with the relative values $h_{K_r,B(a,R)}^*(a)=-1$ on small concentric balls.

\begin{lemma}\label{lem:affine-local-weighted-regularity}
Let $q$ be bounded and upper semicontinuous on $K$, and suppose $\delta>0$ is a global raw-contact parameter for $q$. If $K$ is locally pluriregular at $a$, then $F_{\delta,K,q}(a)=q(a)$ and $K$ is locally $(\delta,q)$-regular there. These equalities persist at every later growth parameter.
\end{lemma}
\begin{proof}
Put $M_r=\sup_{K_r}q$. Raw contact, compact-set monotonicity and the obstacle sandwich give
\[
 q(a)\leq F_{\delta,K,q}(a)
 \leq F_{\delta,K_r,q}(a)
 \leq M_r+\delta V_{K_r}^*(a)=M_r.
\]
Since $q$ is upper semicontinuous and $a\in K_r$, one has $M_r\downarrow q(a)$, giving the global equality.

Every compact subset $E\subseteq K$ inherits raw contact: on $E$ its raw envelope lies between the original envelope $q$ and its obstacle $q$. For fixed $s>0$, apply the same bounds with $r<s$ and $K_s$ in place of $K$, then let $r\downarrow0$. This proves local weighted regularity. The raw contact ray contains every later parameter by \cref{thm:affine-growth-contact-ray}; applying the proved criterion at each of them gives persistence.
\end{proof}

\begin{theorem}\label{thm:affine-local-strict-weight}
Let $a\in K$, $\tau>0$ and $\Psi\in\mathcal L_\tau\cap C(\mathbb C^n)$, and put $q=\Psi|_K$. Suppose that, for some $R,c>0$,
\[
 \ddc\Psi\geq c\omega_0\quad\text{on }B(a,R),
 \qquad\omega_0=\ddc\lVert z\rVert^2.
\]
For $0<r<R$ and $\delta\geq\tau$,
\begin{equation}\label{eq:affine-local-strict-weight-defect}
 F_{\delta,K,q}(a)-\Psi(a)
 \geq cr^2\bigl(1+h_{K_r,B(a,R)}^*(a)\bigr).
\end{equation}
Regular contact at $a$ for one $\delta\geq\tau$ is equivalent to local pluriregularity at $a$. In that case regular contact holds for every $\delta\geq\tau$.
\end{theorem}
\begin{proof}
Set $B=B(a,R)$ and $E=K_r$. For a relative competitor $u$, the function
\[
 H_u(z)=\Psi(z)-c\lVert z-a\rVert^2+cr^2(u(z)+1)
\]
is plurisubharmonic on $B$. Indeed, the current inequality makes its first two terms plurisubharmonic, and the last term is plurisubharmonic. Since $u\leq0$, one has $H_u\leq\Psi$ where $\lVert z-a\rVert\geq r$. Paste $W=\max(\Psi,H_u)$ on $B$ with $W=\Psi$ outside $B$. The pieces agree on an outer collar, so $W$ is globally plurisubharmonic. It belongs to $\mathcal L_\tau$ because it equals $\Psi$ outside the ball. On $E$, $u=-1$ gives $H_u\leq\Psi$; on $K\setminus E$ the distance bound gives the same inequality. Thus $W=q$ everywhere on $K$ and is a raw competitor for each $\delta\geq\tau$. Taking the supremum of $H_u$ near $a$ and regularizing, with the continuous $\Psi$ retained, proves \eqref{eq:affine-local-strict-weight-defect}.

The global extension $\Psi$ itself gives raw contact for every $\delta\geq\tau$. If regular contact holds at one such parameter, the left side of the estimate is zero. Each relative value must then be $-1$. \Cref{lem:relative-affine-zero-value} excludes polar pieces and gives $V_{K_r}^*(a)=0$ for all $r<R$; larger pieces follow by monotonicity. Conversely, local pluriregularity and \cref{lem:affine-local-weighted-regularity} give regular contact at every $\delta\geq\tau$.
\end{proof}
\index{Envelope!locally strict weight}

In this regular case, $V_K^*(a)=0$ by compact-set monotonicity and the initial growth derivative at $\tau$ is zero by persistence and \cref{cor:affine-slope-regular-contact}.

\Needspace{10\baselineskip}
\begin{corollary}\label{cor:affine-weight-degenerate-hessian}\index{Contact set!weight Hessian degeneracy}
Let $\tau>0$, $\Psi\in\mathcal L_\tau\cap C^2(\mathbb C^n)$ and $a\in K$. If $F_{\tau,K,\Psi|_K}(a)=\Psi(a)$ but regular contact fails at some later parameter, then
\[
 (\ddc\Psi)^n(a)=0.
\]
\end{corollary}
\begin{proof}
The smooth Hessian of $\Psi$ is semipositive. A nonzero top power at $a$ would make it positive definite. Continuity would then give $\ddc\Psi\geq c\omega_0$ on a neighborhood of $a$ for some $c>0$. \Cref{thm:affine-local-strict-weight} would force persistence of regular contact, a contradiction.
\end{proof}

